\chapter[2025 October]{October 2025}

\section[2025/10/06]{Thursday, 06 October 2025}
\label{sec:20251006}

I've been extremely bad with the labbook but this is basically what has been happening

Initially the idea was to process the dataset to simplify it substantially decreasing the number of people.

Using coco I did some processing where I could specify the number of people and the natural thought was to use 2
(also as an aside - it has been a while - but I don't think there were many example of 3 or more people so we needed to use 2).
This lead to a problem the model would train but the results were weird see run (d6508c32-9d31-46ae-b8d5-940e9c5f5660) 

So the issue (and concern) was the when doing PCA analysis the results collapsed into a single manifold (a line)
meaning that one component explained almost all the variance. This happened time and time again and we initially 
thought it was a pretty big problem (later I did realise why this was happening because there is only two people it would 
make sense that the variance would be explained by a single component - basically a line between the two people - but anyways).

So the idea was to build a dataset manually (this way I have complete control over the number of people, poses, camera angles etc and
with that have a better idea of what the model is doing).

I did a couple of things not going to go over all the details now but everything from trying to build this myself (but that was a massive
was of time as it's WAY more complicated than I originally thought), using different tools and methods I eventually landed on the following
setup:

I use blender and models from mixamo (owned or built by adobe). These human models come with skeletons and meshes but importantly
poses as well. I then use blender to load random models from mixamo (I've downloaded a few but I'll download a bunch more later)
place them randomly around a scene (facing different directions) run their animations and take photos at different angles (around
a circle). I then map those to a coco format (basically a coco format the only thing I do differently is I add depth into the
annotations - whereas in the proper pipeline I use MiDAS to get depth maps) so a small custom pytorch dataset and datasetloader 
that's almost identical to the "full" pipelines version. Also one other thing that's different is that each image has an
annotation file rather than one large one (that just made the generation code easier but might change that later).

Using this approach I was able to generate 3 daatasets (for now for testing)

\begin{compactitem}
    \item 2 people
    \item 3 people
    \item 4 people
\end{compactitem}

These are the results so far (all the hyper parameters stay the same between the runs):

\renewcommand{\runid}{51d8f9fa-bbc8-44a1-874c-d2edbcb90b3a}

\subsection{Run: \texttt{\runid}}

\todoredefined{Ignore num people required that's an artifact from the dataset processing code I haven't removed}

\lstinputlisting[
    caption={Hyperparameters for run \texttt{\runid}},
    label={code:hyperparameters:\runid}
]{figures/\runid/hyperparameters.json}

\begin{figure}[H]
  \centering
  \includegraphics[width=0.48\linewidth]{figures/\runid/embeddings_epoch_25.png}
  \includegraphics[width=0.48\linewidth]{figures/\runid/embeddings_epoch_50.png}
  \includegraphics[width=0.48\linewidth]{figures/\runid/embeddings_epoch_75.png}
  \includegraphics[width=0.48\linewidth]{figures/\runid/embeddings_epoch_100.png}
  \caption{Embeddings at various epochs for run \texttt{\runid}.}
  \label{fig:embeddings:\runid}
\end{figure}

\begin{figure}[H]
  \centering
  \includegraphics[width=0.9\linewidth]{figures/\runid/training_summary.png}
  \caption{Training summary for run \texttt{\runid}.}
  \label{fig:summary:\runid}
\end{figure}

This is pretty much as expected with 2 people the embeddings collapse into a line again (one component explains almost all the variance).

This is the finally training after 100 epochs:

Epoch 100/100
  Loss: 0.2951
  Pos Distance: 0.8209
  Neg Distance: 0.1236
  
Now onto 3 people 

\renewcommand{\runid}{1ae0a5f5-95e1-42f1-8dc2-5e9ff204ee36}

\subsection{Run: \texttt{\runid}}

\todoredefined{Ignore num people required that's an artifact from the dataset processing code I haven't removed}

\lstinputlisting[
    caption={Hyperparameters for run \texttt{\runid}},
    label={code:hyperparameters:\runid}
]{figures/\runid/hyperparameters.json}

\begin{figure}[H]
  \centering
  \includegraphics[width=0.48\linewidth]{figures/\runid/embeddings_epoch_25.png}
  \includegraphics[width=0.48\linewidth]{figures/\runid/embeddings_epoch_50.png}
  \includegraphics[width=0.48\linewidth]{figures/\runid/embeddings_epoch_75.png}
  \includegraphics[width=0.48\linewidth]{figures/\runid/embeddings_epoch_100.png}
  \caption{Embeddings at various epochs for run \texttt{\runid}.}
  \label{fig:embeddings:\runid}
\end{figure}

\begin{figure}[H]
  \centering
  \includegraphics[width=0.9\linewidth]{figures/\runid/training_summary.png}
  \caption{Training summary for run \texttt{\runid}.}
  \label{fig:summary:\runid}
\end{figure}

Was quite happy with this if you look at keypoint embeddings at epoch 100 (figure~\ref{fig:embeddings:1ae0a5f5-95e1-42f1-8dc2-5e9ff204ee36})
the data is nicely clustered into 3 groups and the variance is more evenly spread across the components.

This is the finally training after 100 epochs:

Epoch 100/100
  Loss: 0.3874
  Pos Distance: 0.8275
  Neg Distance: 0.2080

Then finally for 4 people

\renewcommand{\runid}{2e82e604-a425-4207-8f1b-1468f2743dba}

\subsection{Run: \texttt{\runid}}

\todoredefined{Ignore num people required that's an artifact from the dataset processing code I haven't removed}

\lstinputlisting[
    caption={Hyperparameters for run \texttt{\runid}},
    label={code:hyperparameters:\runid}
]{figures/\runid/hyperparameters.json}

\begin{figure}[H]
  \centering
  \includegraphics[width=0.48\linewidth]{figures/\runid/embeddings_epoch_25.png}
  \includegraphics[width=0.48\linewidth]{figures/\runid/embeddings_epoch_50.png}
  \includegraphics[width=0.48\linewidth]{figures/\runid/embeddings_epoch_75.png}
  \includegraphics[width=0.48\linewidth]{figures/\runid/embeddings_epoch_100.png}
  \caption{Embeddings at various epochs for run \texttt{\runid}.}
  \label{fig:embeddings:\runid}
\end{figure}

\begin{figure}[H]
  \centering
  \includegraphics[width=0.9\linewidth]{figures/\runid/training_summary.png}
  \caption{Training summary for run \texttt{\runid}.}
  \label{fig:summary:\runid}
\end{figure}

This is great (figure~\ref{fig:embeddings:2e82e604-a425-4207-8f1b-1468f2743dba}) looking at the embeddings
at epoch 100 the data is nicely clustered into 4 groups and the variance is evenly spread (around 50 50)

The final results training results look like this

Epoch 100/100
  Loss: 0.3370
  Pos Distance: 0.8023
  Neg Distance: 0.1710
  
\section[2025/10/07]{Friday, 07 October 2025}
\label{sec:20251007}

Okay now onto the results and understanding them

Looking at the results there are a couple takes aways and I'm not sure this is completely correct yet but this is what I'm thinking

Firstly the PCA I think I've covered and that makes sense now which is good.

Secondly a higher (and slower decreasing contrastive loss) when you move from 2 to 3 to 4 people makes sense as well because
the model has to learn to separate more people so the loss should be higher and might just mean that the model is too small (but
this is something that's going to be verified).

Since I'm using cosine similarity, postitives should be high and negatives low (in euclidean the opposite is true).
3 people have the worst seperation at 0.208 (neg distance) which means a higher loss despite having fewer people.
This doesn't really support the initial idea that more people = higher loss so I'm going to assume that 
this particular dataset is just more difficult (maybe the random poses I that were chosen here are more similar (look
alike pair) but I'm not too sure about why that would matter?)

In order to test this lets do another run with 3 people and a larger model

Okay so doing this you get the following results

\renewcommand{\runid}{cd85fc8c-9347-432b-91f6-448a4cb8536b}

\subsection{Run: \texttt{\runid}}

\todoredefined{Ignore num people required that's an artifact from the dataset processing code I haven't removed}

\lstinputlisting[
    caption={Hyperparameters for run \texttt{\runid}},
    label={code:hyperparameters:\runid}
]{figures/\runid/hyperparameters.json}

\begin{figure}[H]
  \centering
  \includegraphics[width=0.48\linewidth]{figures/\runid/embeddings_epoch_25.png}
  \includegraphics[width=0.48\linewidth]{figures/\runid/embeddings_epoch_50.png}
  \includegraphics[width=0.48\linewidth]{figures/\runid/embeddings_epoch_75.png}
  \includegraphics[width=0.48\linewidth]{figures/\runid/embeddings_epoch_100.png}
  \caption{Embeddings at various epochs for run \texttt{\runid}.}
  \label{fig:embeddings:\runid}
\end{figure}

\begin{figure}[H]
  \centering
  \includegraphics[width=0.9\linewidth]{figures/\runid/training_summary.png}
  \caption{Training summary for run \texttt{\runid}.}
  \label{fig:summary:\runid}
\end{figure}

Now this is interesting looking at the embeddings at epoch 100 (figure~\ref{fig:embeddings:cd85fc8c-9347-432b-91f6-448a4cb8536b})
the data isn't nicely clustered but the seperation is better but it's fine at 75?

(I'm going to give all the results here because the neg decreases quite a bit (and then increases again))

Epoch 10/100
  Loss: 0.4092
  Pos Distance: 0.8591
  Neg Distance: 0.3938

Epoch 20/100
  Loss: 0.3644
  Pos Distance: 0.8651
  Neg Distance: 0.2112

Epoch 30/100
  Loss: 0.3366
  Pos Distance: 0.8713
  Neg Distance: 0.1335

Epoch 40/100
  Loss: 0.3203
  Pos Distance: 0.8615
  Neg Distance: 0.0890

Epoch 50/100
  Loss: 0.3017
  Pos Distance: 0.8379
  Neg Distance: 0.0840

Epoch 60/100
  Loss: 0.3003
  Pos Distance: 0.8353
  Neg Distance: 0.0825

Epoch 70/100
  Loss: 0.2930
  Pos Distance: 0.8409
  Neg Distance: 0.0901

Epoch 80/100
  Loss: 0.2841
  Pos Distance: 0.8537
  Neg Distance: 0.1457

Epoch 90/100
  Loss: 0.2872
  Pos Distance: 0.8508
  Neg Distance: 0.1472

Epoch 100/100
  Loss: 0.2939
  Pos Distance: 0.8457
  Neg Distance: 0.1441
  
You can see it gets down to 0.0825 at epoch 60 but then goes back up again to 0.1441 at epoch 100.

This might be because of the number of heads (I set this to be 2 but for consistency I should have left it at 1)

I'm going to test again and see what happens with a "large" model 64 hidden and output dim but with 1 head